\section{Objetivos específicos}

\begin{enumerate}
  \item Investigar y analizar las arquitecturas LSTM y los componentes de los transformadores empleados en el pronóstico de series de tiempo, entre ellos la atención de producto punto escalado, la atención multicabeza y la codificación posicional.
  \item Adaptar e implementar transformadores para regresión con ventanas de entrada $X\in\mathbb{R}^{L\times p}$, codificación temporal y una capa de salida lineal para pronosticar uno o varios pasos.
  \item Diseñar el preprocesamiento y el postprocesamiento de los datos con verificaciones de continuidad, construcción de ventanas, escalado ajustado con datos de entrenamiento y particiones cronológicas que eviten la fuga de información.
  \item Establecer modelos de referencia OLS, Ridge, Lasso, Elastic Net, SVR, Random Forest y perceptrón multicapa bajo condiciones comparables de entrada, entrenamiento y evaluación.
  \item Comparar el desempeño de las redes LSTM y los transformadores con los modelos de referencia mediante RMSE, MAE, MAPE y $R^2$.
  \item Evaluar la predicción a una hora y comparar las estrategias recursiva y de salida múltiple directa para un horizonte de 72 horas.
  \item Analizar el error por paso del horizonte y el costo computacional de los modelos, e identificar oportunidades de mejora.
\end{enumerate}
